\section{Formal game specification}
\label{app:game}

This section completes the game definition in Section~\ref{sec:benac-i}.
The commitment representation, goal satisfaction functions, and terminal
utility follow the action-level negotiation framework of
\citet{benac2026benchmark}. We specify the information available to each
player and the transition rules used in our private-preference setting.

\paragraph{Public instance and private types.}
The public instance specifies the players, commitment coordinates, goal
requirements, scoring modes, proposer schedule, and legal-action rules in
Eq.~\ref{eq:benac-i-game}. A commitment coordinate belongs to a player;
it is not a goal, an item transferred to another player, or an exclusive
choice among alternatives. Different goals may depend on the same
coordinate. Once that coordinate becomes binding, it contributes to every
goal that requires it. A goal's requirements do not indicate which players
want that goal: requirements and preferences are separate parts of the
instance.

The preference assignment $v$ is sampled once and remains fixed throughout
play. Each player observes its own complete vector $v_i$. Any publicly
specified preferences are common knowledge as well. Write
$F=\{(i,g,z)\}$ for these fixed facts and
$\mathcal W=\operatorname{supp}(\mu_0)$ for the admissible joint assignments.
In instances generated from public categorical weights $w_z$, the prior is
obtained by conditioning independent draws of the unfixed entries on the
stated generation constraints:
\begin{equation}
 \mu_0(v)=\frac{1}{Z}
 \mathbf 1[v\text{ agrees with }F]\,
 \mathbf 1[v\text{ satisfies the generation constraints}]
 \prod_{(i,g)\notin\operatorname{dom}(F)} w_{v_{ig}}.
 \label{eq:game-prior}
\end{equation}
Here $Z$ normalizes the distribution. The constraints used in these
instances require each player to want at least one goal and each goal to
have at least one non-neutral player. The weights and fixed facts are
provided with the instance; they are not inferred from the realized world.
Conditioning on these constraints can create dependence between preferences,
so independent background draws do not imply independent posterior beliefs.

\paragraph{State and decision phases.}
A full environment state can be written as
\begin{equation}
 x=(t,C,p,\mathbf q,v,\mathbf r,h),
 \label{eq:game-full-state}
\end{equation}
where $t$ indexes proposal opportunities, $C$ is the binding commitment
matrix, $p$ is either an unresolved offer or $\emptyset$, $\mathbf q$
contains remaining investigation budgets, $\mathbf r=(r^i)_i$ contains
private investigation records, and $h$ is the public event history.
This full state is available to the environment, not to every player.
When $p=\emptyset$ and $t<T$, the actor is $\rho_t$. When an offer is
pending, its recipient acts instead. The index $t$ does not advance until
that offer is accepted or rejected. Thus a proposal and its response count
as two decisions but only one proposal opportunity.

\paragraph{Legal offers.}
At an opportunity of player $i$, an offer is
$o=(j,D_i,D_j)$, where $j\ne i$ and $D_i,D_j$ are sets of new commitment
coordinates owned by the two participants. Let $m$ denote the per-player,
per-offer addition limit specified by $\mathcal L$. Legality requires
\begin{equation}
 D_\ell\subseteq\{k:c_{\ell k}=0\},\qquad
 |D_\ell|\leq m\quad(\ell\in\{i,j\}),\qquad
 |D_i|+|D_j|>0.
 \label{eq:game-offer-legality}
\end{equation}
Any additional forbidden coordinates specified by the instance are excluded.
One participant's addition set may be empty. The limit applies separately
to each participant in this offer; it is not a lifetime commitment limit.
Previously binding coordinates cannot be removed or submitted as new
additions. The native state representation retains those coordinates when
constructing the resulting commitment vectors. The game uses bilateral
offers, not menus of alternative offers.

\paragraph{Transitions.}
When no offer is pending, the scheduled proposer may take one of the
following actions:
\begin{itemize}
 \item \textsc{Pass}: append the public action to $h$, leave $C$ unchanged,
 and advance $t$ by one.
 \item \textsc{Offer}$(j,D_i,D_j)$: append the offer to $h$ and set $p=o$.
 No commitment binds yet and $t$ is unchanged. The recipient must then
 choose \textsc{Accept} or \textsc{Reject}. Acceptance sets
 $c_{\ell k}=1$ for $k\in D_\ell$, $\ell\in\{i,j\}$; rejection leaves
 $C$ unchanged. Either response is public, clears $p$, and advances $t$.
 \item \textsc{Investigate}$(j,g)$: legal only if $j\ne i$ and $q_i>0$.
 The environment appends the public query $(i,j,g)$ to $h$, privately
 appends $(j,g,v_{jg})$ to $r^i$, reduces $q_i$ by one, and advances $t$.
 Commitments are unchanged. Other players' private records are unchanged.
\end{itemize}
Investigation is allowed even when the queried preference is already known
or does not affect the current decision. Its cost is the consumed proposal
opportunity and budget, not a separate utility penalty. It is unavailable
during a pending response. Responses do not consume investigation budgets.

\paragraph{Observations and information sets.}
All players observe the public rules, fixed preference facts, commitments,
proposer schedule, remaining budgets, offers, responses, passes, and query
identities. An investigation answer is revealed only to its investigator.
Player $i$'s view is $I_s^i=(h_s,v_i,r_s^i)$, with the public instance
understood as fixed context. Two histories in the full game belong to the
same information set for $i$ when these views agree. In particular, two
worlds may lead to the same public query while giving its investigator
different private answers. The other players do not receive those answers.

\paragraph{Termination and incentives.}
A full game starts with $C_0=0$, no pending offer, no private query records,
and $q_{i,0}=1$. It ends when all $T$ proposal opportunities have been used,
including the response to an offer made at the final opportunity. Utilities
are then evaluated by Eq.~\ref{eq:benac-i-utility}; there are no intermediate
payoffs. For a binary goal, partial completion yields zero. For a linear
goal, each required coordinate contributes an equal fraction of completion.
An unachieved goal contributes zero even if it is wanted or avoided.
Each player sums its own signed contributions across all goals. Utilities
need not sum to zero, and one player's improvement need not imply either
an improvement or a loss for the others.

\paragraph{Starting a slice.}
A slice records its starting commitments, schedule position, pending offer
if any, remaining budgets, and the observations available at that position.
A state reached by a legal prefix retains the corresponding public history
and each player's own private observations. An explicitly exogenous initial
state is instead a declared starting condition: it does not assert that a
partner previously selected those commitments. Only actual observed choices
provide behavioral evidence under the reference policy. A self-action
specified as an intervention changes the physical state and may deliver a
private answer, but is not itself evidence about a partner's preference.
The reference-policy computation and conditioning conventions are specified
in Appendix~\ref{app:oracle}.
